\section{Zadanie 1}

Niech \(T\) będzie dowolnym drzewem z dwoma wyróżnionymi wierzchołkami: startowym i końcowym. Pionek w sposób losowy porusza się po drzewie rozpoczynając w wierzchołku startowym. W każdym kroku pionek przemieszcza się na losowo wybranego sąsiada aktualnego wierzchołka. Czy dla dowolnego drzewa pionek dojdzie, z prawdopodobieństwem 1, do wierzchołka końcowego?

\subsection*{Rozwiązanie}
Być może da się to zrobić jakoś bardziej elegancko, ale wystarczy tu wykonać trochę szacowania i wyjdzie. Niech \(T\) - drzewo z zadania o wierzchołku końcowym \(t\). Dodatkowo:

\begin{align*}
max\_dist(G) &= max_{v \in V(G)}dist(v, t)\\
\Delta(G) &= max_{v \in V(G)}deg(v)
\end{align*}

...Albo po prostu maksymalna odległość od końcowego i maksymalny stopień w grafie.

Dla naszego grafu, dla każdego wierzchołka innego niż \(t\) przypisujemy liczby do każdej wychodzącej z niego krawędzi, zaczynając od \(1\) dla tej krawędzi, która prowadzi do \(t\). Każdy wierzchołek ma maksymalnie \(\Delta(G)\) krawędzi wychodzących oraz jest o maksymalnie \(max\_dist(G)\) krawędzi oddalony od \(t\).

Oznacza to, że losując kolejno krawędzie, którymi się poruszamy, wylosowanie ciągu \(max\_dist(G)\) jedynek doprowadzi nas do \(t\) niezależnie od wierzchołka, z którego rozpoczęliśmy losowanie.

Losujemy więc liczby ze zbioru \([\Delta(G)]\) w blokach po \(max\_dist(G)\) i poruszamy się wylosowanymi krawędziami, aż nie dotrzemy do \(t\) (jeśli wylosujemy liczbę większą niż liczba dostępnych krawędzi, ignorujemy rzut, tak jakby rzucamy ponownie).

Prawdopodobieństwo nie wystąpienia bloku \(max\_dist(G)\) jedynek po wyrzuceniu \(k\) kolejnych bloków prezentuje się następująco:

\begin{align*}
1 - \left(1 - \frac 1 {\Delta(G)^{max\_dist(G)}}\right)^k
\end{align*}

Niech \(B_k\) to zdarzenie polegające odniesieniu sukcesu nie później niż \(k\)-ty rzut. Ograniczamy zdarzenie z zadania przez zdarzenie \(\bigcup_{k=0}^\infty B_k\) i przez \(B_i \subseteq B_{i+1}\) dostajemy:

\[
P(\bigcup_{k=0}^\infty B_k) = \lim_{k \to \infty}P(B_k) = \lim_{k \to \infty}\left[1 - \left(1 - \frac 1 {\Delta(G)^{max\_dist(G)}}\right)^k\right] = 1
\]

Stąd mamy prawdopodobieństwo dotarcia pionka do wierzchołka końcowego równe 1.

\subsection*{Odpowiedź}
\[
\boxed{\text{Tak, dla dowolnego drzewa pionek dojdzie do wierzchołka końcowego.}}
\]
